\section{容器与环境：只给 Chatbot 输入框是不负责任的}

上一章讲 token，本章讲承载 token 的产品形态。明超平认为，“壳”这个词低估了应用公司的价值，更准确的说法是容器和环境。一个 Chatbot 输入框只是入口，它把用户意图交给模型；而环境会影响用户如何表达、如何行动、如何反馈、如何分享作品。环境是人的反应器，也决定模型能力如何被组织起来。

\begin{figure}[H]
\centering
\includegraphics[width=0.94\textwidth]{container-environment.png}
\caption{壳、容器与环境：只给 Chatbot 输入框不够，环境是人的反应器。自制概念图，依据 01:16:33--01:21:18 对谈内容整理。}
\end{figure}

\begin{knowledgebox}{读图：Shell 只是入口，Environment 才能塑造行为}
左边 Shell 只是承接输入输出；右边 Environment 承载上下文、行动、反馈和人的反应。YouWare 的产品假设是：AI 创作需要一个能发布、互动、分享和继续演化的环境，而不是把代码扔给用户自己部署。
\end{knowledgebox}

\subsection{Coding 像 Camera：新工具创造新人群}

本节解释 YouWare 为什么押 coding 创作。明超平把 coding 类比 camera：早期谈相机，说的是拿单反的人；手机摄影出现后，出现了“手机摄影师”这个新人群。Coding 也是如此。过去 coding 属于程序员，AI coding 之后会出现 Vibe Coder：不一定懂完整工程体系，但能用模型、审美和意图生成作品的人。

\begin{figure}[H]
\centering
\includegraphics[width=0.96\textwidth]{vibe-coder-new-crowd.png}
\caption{Vibe Coder 新人群：Coding 像 Camera，工具变化会创造新创作者人群。自制概念图，依据 01:21:18--01:37:12 对谈内容整理。}
\end{figure}

\begin{knowledgebox}{读图：新人群是早期趋势变量}
图中 Camera 从单反用户转向手机摄影师，Coding 从工程师转向 Vibe Coder。关键不是原有专业人群效率提升多少，而是新工具是否创造了更大的新人群和新内容形态。
\end{knowledgebox}

\begin{importantbox}{YouWare 的最小切口}
明超平第一次写 YouWare，是为了解决“AI 写出来的 HTML/网站/小游戏无法被普通人分享和使用”的问题。把代码贴进去，得到一个可访问的网站，这个切口把 coding 从本地工程环境拉进创作和分发环境。
\end{importantbox}

\subsection{本章小结}

容器和环境是 AI 应用公司的核心价值。模型会变强，但用户需要一个能承接创作、分享、互动和反馈的场。YouWare 试图把 coding 从工程任务变成新的创作方式。

